\section{Does what gets verified matter later?}\label{sec:downstream}

Allocation only matters if the record the budget failed to recover changes what
the agent later does. We test that directly by randomising it. A second agent
inherits the same situation together with a carried set of provenance records,
and which records are carried is assigned at random rather than chosen by the
first agent. This breaks the dependence between what an agent wanted to check
and what it ends up holding.

Carrying the source record for the corrupted memory moves the rate at which the
later decision goes in the corrupted direction from \ThreeBOthK/\ThreeBOthN{}
to \ThreeBTgtK/\ThreeBTgtN{}, and unguarded commitment --- committing without
naming the constraint --- from \ThreeBNGOth{} to \ThreeBNGTgt{}
($\ThreeBCMH$, all \ThreeBAllModels{} models in the same direction). A randomised-instrument replay --- re-running only the second decision as a
fresh session, with recovered records reproduced from each stored episode ---
agrees with the direct estimate within a point; we report it as convergent
evidence only, since the exclusion restriction it needs is untestable here.

This identifies the \emph{downstream} half of the chain: what provenance the
agent holds causally changes the decision it reaches. Together with
\S\ref{sec:directs}, which identifies the upstream half, the two halves are
established separately. We do not claim an identified end-to-end mediation from
working plan through verification allocation to later harm; no single design
here carries that, and combining two separately identified halves is not the
same as identifying the chain. Nor do we present this as the novel mechanism:
inserting the decisive contradicting record into an agent's context and
watching the decision change is close to what one would expect. Its role is to
establish that the allocation studied upstream has operational consequence.
